\documentclass[aspectratio=169]{beamer}
\usetheme{Madrid}
\usepackage[spanish,es-nodecimaldot]{babel}
\usepackage[utf8]{inputenc}
\usepackage[T1]{fontenc}
\usepackage{amsmath}
\usepackage{amssymb}
\usepackage{graphicx}
\usepackage{booktabs}
\usepackage{xcolor}

% Color UNS
\definecolor{azulUNS}{RGB}{30,60,120}
\usecolortheme[named=azulUNS]{structure}

\title[MAD1 -- U6]{Métodos de Análisis de Datos I}
\subtitle{Unidad 6 -- Bloque 4: El data scientist en contexto}
\author{Departamento de Matemática -- UNS}
\date{2026}

\begin{document}

\begin{frame}
  \titlepage
\end{frame}

\begin{frame}{Contenidos}
  \tableofcontents
\end{frame}

%-----------------------------------------------
\section{Más allá del modelo}
%-----------------------------------------------

\begin{frame}{El trabajo real de un analista de datos}
  En los cursos aprendemos a construir modelos. En el trabajo real, construir el modelo es solo una parte.

  \bigskip
  \begin{columns}
    \column{0.5\textwidth}
      \textbf{Lo que se enseña:}
      \begin{itemize}
        \item Preprocesamiento
        \item Selección de modelo
        \item Entrenamiento y evaluación
        \item Métricas
      \end{itemize}
    \column{0.5\textwidth}
      \textbf{Lo que también ocupa tiempo real:}
      \begin{itemize}
        \item Entender el problema de negocio
        \item Conseguir y limpiar los datos
        \item Negociar qué pregunta es respondible
        \item Comunicar resultados
        \item Mantener y monitorear modelos en producción
      \end{itemize}
  \end{columns}

  \bigskip
  \begin{alertblock}{Estimación práctica}
    En proyectos reales, entre el 60\% y el 80\% del tiempo se gasta en datos: obtenerlos, entenderlos, limpiarlos, validarlos. El modelado ocupa una fracción menor de lo que se imagina.
  \end{alertblock}
\end{frame}

%-----------------------------------------------
\section{Ética y sesgo en los datos}
%-----------------------------------------------

\begin{frame}{Sesgo: el problema empieza en los datos}
  Un modelo aprende de los datos. Si los datos tienen sesgos, el modelo los amplifica.

  \bigskip
  \textbf{Tipos de sesgo frecuentes:}
  \begin{itemize}
    \item \textbf{Sesgo de selección:} los datos no representan a la población objetivo. \\
    \textit{Ejemplo: un modelo de crédito entrenado solo con clientes que accedieron al crédito, sin incluir a los rechazados.}
    \item \textbf{Sesgo histórico:} los datos reflejan decisiones pasadas injustas. \\
    \textit{Ejemplo: contratación histórica con discriminación de género $\rightarrow$ modelo que perpetúa esa discriminación.}
    \item \textbf{Sesgo de medición:} el instrumento que genera los datos no mide lo mismo para todos los grupos. \\
    \textit{Ejemplo: detección de enfermedades con herramientas calibradas para un solo grupo demográfico.}
  \end{itemize}
\end{frame}

\begin{frame}{Sesgo: casos reales}
  \begin{columns}
    \column{0.5\textwidth}
      \textbf{COMPAS (2016):}\\
      Sistema de predicción de reincidencia criminal usado en EE.UU.\\
      Un análisis periodístico mostró que asignaba puntuaciones de riesgo más altas a acusados afroamericanos que a blancos con perfiles similares.\\
      \bigskip
      \textbf{Amazon Recruiting (2018):}\\
      Amazon descartó un sistema de selección de CVs porque penalizaba sistemáticamente a mujeres, habiendo aprendido de contrataciones históricas dominadas por hombres.
    \column{0.5\textwidth}
      \textbf{Reconocimiento facial:}\\
      Estudios (MIT Media Lab, 2018) mostraron que sistemas comerciales tenían tasas de error hasta 34 puntos porcentuales más altas para mujeres de piel oscura que para hombres de piel clara.\\
      \bigskip
      \begin{alertblock}{La lección}
        La objetividad del algoritmo es una ilusión si los datos de entrenamiento codifican injusticia histórica.
      \end{alertblock}
  \end{columns}
\end{frame}

\begin{frame}{Equidad algorítmica: no hay una sola definición}
  Intuitivamente, un modelo ``justo'' trata igual a todos. Formalmente, hay varias definiciones incompatibles entre sí:

  \bigskip
  \begin{itemize}
    \item \textbf{Paridad demográfica:} igual tasa de positivos en todos los grupos.
    \item \textbf{Igualdad de oportunidades:} igual tasa de verdaderos positivos (TPR) entre grupos.
    \item \textbf{Calibración:} igual significado de la puntuación de riesgo entre grupos.
  \end{itemize}

  \bigskip
  \begin{alertblock}{Resultado teórico (Chouldechova, 2017)}
    Cuando las tasas base de la variable objetivo difieren entre grupos, \textbf{no es posible} satisfacer simultáneamente paridad demográfica, igualdad de oportunidades y calibración.\\
    Elegir una definición de equidad implica sacrificar otra.
  \end{alertblock}
\end{frame}

%-----------------------------------------------
\section{Cuándo NO usar ML}
%-----------------------------------------------

\begin{frame}{Cuándo NO usar machine learning}
  ML no es la solución a todo. Hay situaciones donde \textbf{no corresponde} o \textbf{no ayuda}.

  \bigskip
  \begin{columns}
    \column{0.5\textwidth}
      \textbf{No usar ML cuando:}
      \begin{itemize}
        \item El problema se resuelve con una regla simple y conocida.
        \item No hay suficientes datos.
        \item El proceso subyacente cambia tan rápido que cualquier modelo queda desactualizado.
        \item La explicabilidad es un requisito no negociable y el modelo no la provee.
        \item El costo de un error es muy alto y no hay forma de auditarlo.
      \end{itemize}
    \column{0.5\textwidth}
      \begin{exampleblock}{Ejemplo}
        Una empresa quiere predecir si un cliente pagará a tiempo.\\
        Tiene 80 registros históricos.\\
        \bigskip
        Un modelo de ML no tiene suficientes datos para generalizar. Una regla heurística (``clientes con más de 3 atrasos previos $\rightarrow$ riesgo alto'') puede ser más robusta y explicable.
      \end{exampleblock}
  \end{columns}
\end{frame}

\begin{frame}{La maldición de la complejidad innecesaria}
  Más complejo $\neq$ mejor.

  \bigskip
  \begin{columns}
    \column{0.55\textwidth}
      \textbf{Principio de parsimonia:}\\
      Entre dos modelos con rendimiento similar, preferir el más simple. Razones:
      \begin{itemize}
        \item Más fácil de interpretar y auditar.
        \item Menos propenso a sobreajuste en datos pequeños.
        \item Más fácil de mantener en producción.
        \item Menos costoso computacionalmente.
        \item Más fácil de explicar a quienes toman decisiones.
      \end{itemize}
    \column{0.45\textwidth}
      \begin{alertblock}{Regla práctica}
        Antes de usar XGBoost o una red neuronal, probar regresión logística o un árbol de decisión simple.\\
        Si la ganancia de precisión no justifica la pérdida de interpretabilidad, quedarse con el modelo simple.
      \end{alertblock}
  \end{columns}
\end{frame}

%-----------------------------------------------
\section{Limitaciones del modelo vs. limitaciones del dato}
%-----------------------------------------------

\begin{frame}{Distinguir: ¿falla el modelo o fallan los datos?}
  Cuando un modelo tiene mal desempeño, hay dos causas posibles muy distintas:

  \bigskip
  \begin{columns}
    \column{0.5\textwidth}
      \textbf{Limitación del modelo:}
      \begin{itemize}
        \item El algoritmo es inadecuado para la estructura del problema.
        \item Hiperparámetros mal ajustados.
        \item Sobreajuste o subajuste.
      \end{itemize}
      \bigskip
      \textbf{Solución:} cambiar el modelo, ajustar, regularizar.
    \column{0.5\textwidth}
      \textbf{Limitación del dato:}
      \begin{itemize}
        \item La variable objetivo está mal definida o mal medida.
        \item Las variables predictoras no capturan el fenómeno real.
        \item Los datos son insuficientes, sesgados o desactualizados.
      \end{itemize}
      \bigskip
      \textbf{Solución:} mejorar la recolección, redefinir el problema, recolectar más datos.
  \end{columns}

  \bigskip
  \begin{alertblock}{Error frecuente}
    Cambiar el modelo cuando el problema es el dato. Ningún algoritmo puede extraer información que no está en los datos.
  \end{alertblock}
\end{frame}

%-----------------------------------------------
\section{El rol del analista como traductor}
%-----------------------------------------------

\begin{frame}{El analista como traductor}
  El analista de datos ocupa una posición única: entiende tanto el lenguaje de los datos como el del dominio.

  \bigskip
  \begin{columns}
    \column{0.5\textwidth}
      \textbf{Traducción técnica $\rightarrow$ dominio:}
      \begin{itemize}
        \item ``AUC = 0.85'' $\rightarrow$ ``Detectamos 8 de cada 10 casos relevantes.''
        \item ``Coeficiente $\beta = 0.3$'' $\rightarrow$ ``Un año más de antigüedad se asocia con un 30\% más de probabilidad de retención.''
        \item ``Cluster 3'' $\rightarrow$ ``Clientes jóvenes de alto valor con compras nocturnas.''
      \end{itemize}
    \column{0.5\textwidth}
      \textbf{Traducción dominio $\rightarrow$ técnica:}
      \begin{itemize}
        \item ``Queremos predecir quién va a comprar de nuevo'' $\rightarrow$ ¿clasificación, regresión, ranking?
        \item ``Necesitamos detectar fraude'' $\rightarrow$ ¿cuál es el costo relativo de falsos positivos vs. falsos negativos?
        \item ``¿Por qué el modelo rechazó este caso?'' $\rightarrow$ XAI local.
      \end{itemize}
  \end{columns}
\end{frame}

\begin{frame}{Preguntar bien antes de modelar}
  Antes de abrir un notebook, hay preguntas que deben responderse:

  \bigskip
  \begin{enumerate}
    \item \textbf{¿Cuál es la decisión que se tomará con este análisis?} \\
    No ``queremos entender los datos'', sino ``si el modelo dice X, haremos Y''.
    \item \textbf{¿Cuál es el costo de equivocarse?} \\
    Falso positivo vs. falso negativo: ¿cuál es peor en este contexto?
    \item \textbf{¿Los datos disponibles son adecuados para la pregunta?} \\
    Antes de modelar, no después.
    \item \textbf{¿Existe ya una solución simple que funcione?} \\
    Si hay una regla de negocio que funciona bien, ¿por qué agregar complejidad?
    \item \textbf{¿Cómo se va a usar el modelo en producción?} \\
    Un modelo que nadie implementa no tiene valor.
  \end{enumerate}
\end{frame}

%-----------------------------------------------
\section{El ciclo de vida de un proyecto de datos}
%-----------------------------------------------

\begin{frame}{El ciclo de vida de un proyecto de datos}
  \begin{columns}
    \column{0.5\textwidth}
      \textbf{Fases típicas:}
      \begin{enumerate}
        \item Definición del problema
        \item Recolección y exploración de datos
        \item Preprocesamiento y feature engineering
        \item Modelado y evaluación
        \item Interpretación y comunicación
        \item Despliegue (deployment)
        \item Monitoreo y mantenimiento
      \end{enumerate}
    \column{0.5\textwidth}
      \begin{exampleblock}{El ciclo no es lineal}
        En la práctica, se vuelve atrás constantemente.\\
        La exploración revela que la pregunta estaba mal formulada.\\
        El modelo revela que los datos tienen errores.\\
        El despliegue revela que el contexto cambió.\\
        \bigskip
        Iterar no es fracasar: es el proceso normal.
      \end{exampleblock}
  \end{columns}
\end{frame}

\begin{frame}{Data drift: cuando el mundo cambia}
  Un modelo entrenado hoy puede degradarse con el tiempo sin que nadie lo cambie.

  \bigskip
  \textbf{¿Por qué?}
  \begin{itemize}
    \item \textbf{Data drift:} la distribución de las variables de entrada cambia. \\
    \textit{Ejemplo: un modelo de precios entrenado antes de una crisis inflacionaria.}
    \item \textbf{Concept drift:} la relación entre variables y objetivo cambia. \\
    \textit{Ejemplo: patrones de fraude cambian cuando los defraudadores adaptan su comportamiento.}
  \end{itemize}

  \bigskip
  \textbf{Consecuencia:} los modelos en producción deben monitorearse. Las métricas de un modelo no son permanentes.

  \bigskip
  \begin{alertblock}{Regla práctica}
    Antes de desplegar un modelo, definir cómo y cuándo se va a reevaluar. Un modelo sin monitoreo es una deuda técnica.
  \end{alertblock}
\end{frame}

%-----------------------------------------------
\section{Cierre}
%-----------------------------------------------

\begin{frame}{Ideas para llevarse}
  \begin{enumerate}
    \item El trabajo real de un analista incluye mucho más que entrenar modelos.
    \item Los sesgos en los datos se amplifican en los modelos: la ``objetividad'' del algoritmo es una ilusión.
    \item No existe una definición única de equidad algorítmica: elegir una implica sacrificar otra.
    \item Antes de usar ML, preguntar si una solución más simple alcanza.
    \item Distinguir limitaciones del modelo de limitaciones del dato: son problemas distintos con soluciones distintas.
    \item El analista traduce entre el lenguaje técnico y el lenguaje del dominio: en ambas direcciones.
    \item Los modelos en producción se degradan: monitorear no es opcional.
  \end{enumerate}
\end{frame}

\end{document}
